\documentclass[11pt]{amsart}
\usepackage{geometry}                % See geometry.pdf to learn the layout options. There are lots.
\geometry{letterpaper}                   % ... or a4paper or a5paper or ... 
%\geometry{landscape}                % Activate for for rotated page geometry
%\usepackage[parfill]{parskip}    % Activate to begin paragraphs with an empty line rather than an indent
\usepackage{graphicx}
\usepackage{amssymb}
\usepackage{epstopdf}
\DeclareGraphicsRule{.tif}{png}{.png}{`convert #1 `dirname #1`/`basename #1 .tif`.png}

\title{EJEMPLO 4.4}
%\date{}                                           % Activate to display a given date or no date

\begin{document}
\maketitle
%\section{}
%\subsection{}
\noindent Consideremos el arreglo de enteros del ejemplo anterior. Suponiendo que el límite inferior es igual a -13 y el límite superior igual a 56, su representación queda como se muestra en la figura 4.4.\\

\begin{table}[htbp]
\centering
\label{tab:corrida}
\resizebox{5cm}{!} {
\begin{tabular}{|c|c|c|c|c|c|c|c|}
\hline
 &  &  &  &  .....  & &   \\
 \hline
 \end{tabular}
 }
 \end{table}
 
 \begin{table}[htbp]
\centering
\label{tab:corrida}
\resizebox{5cm}{!} {
\begin{tabular}{cccccccc}
-13 & -12 & -11 &   & &  &56\\
\end{tabular}
}
 \end{table}

\begin{itemize}
  \item $NTE = (56 - (-13) + 1) - 70$
  \item \noindent Cada elemento del arreglo será un número entero y podrá accesarse por medio de un índice que será un valor comprendido entre -13 y 56. \\
\end{itemize} 

\noindent Así por ejemplo: \\

\noindent ARRE[-13] hace referencia al primer elemento del arreglo.\\
                ARRE[-12] hace referencia al segundo elemento del arreglo.\\
                .....\\
                ARRE[1] hace referencia al decimoquinto elemento del arreglo.\\
                .....\\
                ARRE[56] hace referencia al último elemento del arreglo.\\
                
\noindent \textbf {Operaciones con arreglos} \\

\noindent  A continuación se presentan las operaciones más comunes en arreglos:\\
\begin{itemize}
  \item Lectura/Escritura.
  \item Asignación.
  \item Actualización: Inserción,
                                 Eliminación,
                                 Modificación.
  \item Ordenación.
  \item Búsqueda.
  \end{itemize}
\noindent Como los arreglos son datos estructurados, muchas de estas operaciones no pueden llevarse a cabo de manera global, sino que se debe trabajar sobre cada elemento. A continuación se analizará cada una de estas operaciones. Ordenación y búsqueda, serán analizadas en los problemas que presentaremos posteriormente.\\

\noindent \textbf {Explicación de las variables} \\

\noindent \quad N: Variable de tipo entero. Almacena el número actual de elementos del arreglo.\\

\noindent \quad X: Variable de tipo entero. Representa al valor que se va a eliminar.\\

\noindent \quad I: Variable de tipo entero. Se utiliza como variable de control del ciclo externo y como índice del arreglo V.\\

\noindent \quad BAND: Variable de tipo booleano. Se inicializa en falso. Cambia su valora verdadero si se encuentra el elemento a eliminar en el arreglo, en cuyo caso se interrumpe el ciclo.\\

\noindent \quad K: Variable de tipo entero. Se utiliza como variable de control del ciclo interno y como índice del arreglo A.(‘F’).\\

\noindent \quad A: Arreglo unidimensional de tipo entero. Su capacidad máxima es de 100 elementos. \\
\end{document}  
